El sistema de infusión opera bajo una arquitectura de lazo cerrado compuesta por ocho modelos atómicos DEVS interconectados. La simulación es autocontenida: todos los estímulos son generados por componentes internos del modelo y la única salida observable desde el exterior es la notificación de alarmas al entorno hospitalario.

\section{Componentes del Sistema}
El modelo acoplado se compone de los siguientes elementos:
\begin{itemize}
    \item \textbf{Generador de Órdenes Médicas ($G_{om}$):} Simula la llegada estocástica de prescripciones que establecen el caudal objetivo de infusión (entre 0 y 200 ml/h).
    \item \textbf{Generador de Confirmaciones del Enfermero ($G_{ce}$):} Simula la intervención del personal de enfermería, emitiendo confirmaciones a intervalos aleatorios.
    \item \textbf{Generador de Fin de Bolsa ($G_{fb}$):} Monitorea el volumen restante de líquido en la bolsa de infusión y emite una alerta preventiva cuando estima que quedan 60 segundos para su agotamiento.
    \item \textbf{Controlador de Bomba ($C$):} Núcleo lógico del sistema. Recibe las órdenes médicas, las lecturas del sensor, las alertas de la bolsa y las confirmaciones del enfermero. Evalúa desviaciones de caudal y plazos de seguridad para coordinar las acciones del actuador y la activación jerárquica de alarmas.
    \item \textbf{Actuador ($A$):} Modela la respuesta física de la bomba, aplicando una latencia mecánica aleatoria a las órdenes del controlador.
    \item \textbf{Sensor de Flujo ($G_{sf}$):} Mide el caudal real entregado por el actuador, inyectando un ruido Gaussiano que simula las imprecisiones de un instrumento físico.
    \item \textbf{Módulo de Alarmas ($M_a$):} Gestiona la presentación de alarmas al entorno hospitalario, implementando un esquema jerárquico de prioridades y un ciclo de re-notificación para alarmas críticas.
    \item \textbf{Registrador de Eventos ($R_e$):} Componente sumidero que almacena la traza de auditoría de todas las decisiones tomadas por el controlador.
\end{itemize}

\section{Señales Internas}
El Controlador de Bomba ($C$) constituye el punto de convergencia de las señales del sistema. Recibe cuatro tipos de estímulos internos: prescripciones médicas ($G_{om}$), lecturas del sensor de flujo ($G_{sf}$), alertas de agotamiento de la bolsa ($G_{fb}$) y confirmaciones del enfermero ($G_{ce}$). Como respuesta, emite comandos físicos hacia el Actuador (\texttt{ajustarCaudal}, \texttt{detenerBomba}), despacha alertas de tres niveles de severidad (baja, media y crítica) hacia el Módulo de Alarmas ($M_a$), y documenta cada decisión mediante tokens de auditoría enviados al Registrador de Eventos ($R_e$).

\section{Restricciones y Requisitos Temporales}
Para garantizar la seguridad del paciente, el diseño satisface los siguientes parámetros críticos:
\begin{itemize}
    \item \textbf{Latencia del actuador:} La bomba responde a cada orden del controlador con un retardo mecánico aleatorio de entre 0 y 0.5 segundos.
    \item \textbf{Frecuencia de muestreo:} El sensor de flujo emite lecturas cada $1$ segundo.
    \item \textbf{Detección de desvíos:} Un desvío de caudal mayor al $10\%$ respecto al objetivo es tolerado hasta 5 segundos, instante en que se emite una \textbf{alarma media}. Si persiste hasta los 10 segundos, se escala a \textbf{alarma crítica} y se detiene la bomba.
    \item \textbf{Ciclo de alarmas críticas:} Requieren confirmación humana en $30$ segundos; de no recibirla, la alarma se repite cada $10$ segundos indefinidamente.
    \item \textbf{Fin de bolsa:} Tras detectarse que quedan 60 segundos de líquido, la bomba opera ese tiempo adicional y luego se detiene automáticamente.
\end{itemize}
